\documentclass[a4paper]{article}
\usepackage[pdftex]{graphicx}
\usepackage[utf8]{inputenc}
\usepackage{enumerate}
\usepackage{amssymb}
\usepackage{colortbl}
\usepackage{icomma}
\usepackage{siunitx}
\sisetup{locale=DE}
\usepackage{geometry}
\geometry{a4paper, top=15mm, left=15mm, right=15mm, bottom=15mm, 
	headsep=10mm, footskip=12mm}
\usepackage{href-ul}
\hypersetup{ 
	colorlinks=true, 
	linkcolor=blue, 
	urlcolor=blue}
\begin{document}
	
{\LARGE \bf $ \mathbf{F= m \cdot g} $ or "a Journey to the Stars"}
\vspace{1,5 cm}

\noindent	
	{\bf Imagine: It's a normal school day.}
	You're standing in the schoolyard at recess Suddenly, a silver spaceship lands silently among the trees. A ramp extends, and three friendly, smiling aliens step out.
	"Come with us!" they say. "We're going on a research trip through the solar system – and you can come along!"
	Of course, you say yes! But they explain:
	"You're only allowed to take seven things with you. And before you board, you have to weigh everything – we want to know how heavy your luggage is on Earth."
	
	\begin{enumerate}[1.]
		\item Divide these items of luggage into useful, somewhat useful, and not useful:\\
		\fbox{Vacuum cleaner {\bf K}}
		\noindent\fbox{Spacesuit with life support system {\bf J}}
		\fbox{Stuffed toy {\bf Z}}
		\fbox{Toothbrush {\bf E}}\\
		\fbox{Suitcase full of euros {\bf M}}
		\fbox{Seed packets with earth flowers {\bf R}}
		\fbox{Notebook and pen {\bf I}}\\
		\fbox{Small globe as a guest gift {\bf P}}
		%\fbox{Physics 7 textbook {\bf I}}
		\fbox{Digital camera {\bf T}}
		\fbox{Piano {\bf L}}
		%\fbox{Chocolate {\bf T}}
		%\fbox{Favorite book {\bf }}
		%\fbox{Travel first-aid kit {\bf }}
		\fbox{Umbrella {\bf W}}
		\fbox{Clothing {\bf U}}
		\fbox{Cigarettes {\bf N}}\\
		\fbox{E-scooter {\bf O}}
		%\fbox{Tablet with offline star map {\bf U}}
		\fbox{Cat {\bf S}}\\
		
		So you place the objects on a scale: spacesuit, notebook, digital camera...
		You carefully write down the numbers – in kilograms. Use two valid digits for all calculations!
		
		\item Sort the seven useful items by mass. Start with the heaviest! The letters form a solution word.
		\item What is the weight of these items on Earth? The acceleration due to gravity on Earth is $g_{Earth}=\SI{9.81}{\newton\per\kilogram}$
		
		\renewcommand{\arraystretch}{2}
		\begin{tabular}{|p{3 cm}|p{1.5 cm}|p{1.5 cm}|p{1.5 cm}|p{1.5 cm}|p{1.5 cm}|p{1.5 cm}|p{1.5 cm}|}
			\hline
			Piece of luggage & & & & & & & \\
			\hline
			Mass $m$ & $\SI{80}{\kilogram}$ & & $\SI{0.75}{\kilogram}$& & $\SI{200}{\gram}$&& $\SI{5.0}{\gram}$\\
			\hline
			Gravity on Earth $F_{g, Earth}$ & & $\SI{12}{\newton}$ & & $\SI{1.0}{\newton}$ & & $\SI{0.20}{\newton}$ & \\
			\hline
		\end{tabular}
		
		
		
		\fbox{$\SI{20}{\gram}$}
		\fbox{$\SI{0.10}{\kilogram}$}
		\fbox{$\SI{1.2}{\kilogram}$}
		\fbox{$\SI{2.0}{\newton}$}
		\fbox{$\SI{0.78}{\kilo\newton}$}
		\fbox{$\SI{49}{\milli\newton}$}
		\fbox{$\SI{7.4}{\newton}$}
		
		\item What is the mass of these things on the moon?
		
		{\bf Then it starts.} The ramp closes, the spaceship shakes – and suddenly you feel something pushing you into your seat:
		The spaceship accelerates at 7 g!
		You feel seven times heavier than usual.
		Your 40 kg body suddenly "weighs" as much as if you weighed 280 kg. Even your pen feels like lead.
		"Don't worry," says the pilot. "It will only take a few seconds!"
		
		\item What is the spaceship's speed after 100 seconds of acceleration at 7g?
		
		Shortly afterward, everything becomes quiet – weightlessness!
		You float through the cabin, your globe dances beside you, and your notebook floats through the air.
		Everything still has the same mass, but no longer has any gravitational force because you and your belongings are currently falling freely – along with the spaceship.
		
		\item Mars is 86.4 million kilometers from Earth. What cruising speed must the spaceship travel at for the flight there to take 24 hours?
		Please specify the speed in $\SI{}{\kilo\meter\per\hour}$, $\SI{}{\kilo\meter\per\second}$, and $\SI{}{\meter\per\second}$.
		\item How long would the acceleration phase at 7g then last? Specify the time in seconds, minutes, and hours.
		
		After a few hours, Mars comes into view. The spacecraft decelerates gently and lands in a red cloud of dust.
		As you step out, you suddenly feel as light as a balloon—because on Mars, gravity is only about a third as strong as on Earth.
		A backpack that weighs 9 kg on Earth feels as if it only has a weight of 3 kg here.
		
		\item On Mars, the location factor $g_{Mars}=\SI{3.7}{m\per s^2} .$ applies. What is the weight of the spacesuit there?
		You take a selfie of yourself in front of the Martian volcano Olympus Mons and radio it to your parents on Earth. How long will it take to get there if the speed of light is $c=\SI{3e8}{\meter\per\second}$?
		\item How would a Mars-Earth telephone conversation proceed under these conditions?
		
	The aliens smile:
	"Well, Earthling – we have one more stop! We're flying to an exoplanet called Gravimurp-7 near the star Alpha Centauri. We want you to determine the location factor g there."
	You are given a string pendulum and your familiar equipment.
	
	\item You land on the planet "Gravimurp-7." You notice: Your pendulum swings more slowly than on Earth. You attach your digital camera to the force gauge for the gravitational force, and it displays $F_g=\SI{1,2}{\newton}$.
	What is the location factor of the exoplanet?
	
	After a few measurements and calculations, you estimate:
	"Here, the acceleration due to gravity is approximately $\SI{6}{\meter\per\second^2 }$—so slightly weaker than on Earth, but stronger than on Mars!"
	You scatter the bag of earthflower seeds on the damp ground of Gravimurp-7.
	The aliens nod in satisfaction.
	
	"Very good! You've understood that mass always remains the same, but the force of gravity depends on the position factor g."
	Then they gently beam you back to your schoolyard.
	Your belongings are still where you left them – only the speck of dust from Mars on your shoe reminds you
	that you were really there.
	
	\item At home, you sit down at your computer and create an Excel spreadsheet with the mass of everyday objects and their weight on Earth,
	during acceleration in the spaceship, on Mars, and on the exoplanet Gravimurp-7:
	
	\renewcommand{\arraystretch}{2}
	\begin{tabular}{|p{1.5 cm}|p{1.5 cm}||p{3 cm}|p{2 cm}|p{2 cm}|p{2 cm}|p{2 cm}|}
		\hline
		Location factor in g & Location factor in $\SI{}{\meter\per\second^2 }$& & Water bottle 1 kg & Cat 4 kg & Human 70 kg & Smartphone 0.2 kg \\
		\hline
		1 g & $\SI{9.81}{\meter\per\second^2 }$& $F_{g}$ on Earth & & & &\\
		\hline
		7g & $\SI{69}{\meter\per\second^2 }$ & $F_{g}$ in rocket & & & &\\
		\hline
		0.3g & $\SI{3.7}{\meter\per\second^2 }$ & $F_{g}$ on Mars & & & &\\
		\hline
		0.6g & $\SI{6}{\meter\per\second^2 }$ & $F_{g}$ on Gravimurp-7 & & & &\\
		\hline
		
	\end{tabular}
	\end{enumerate}
	
	\fbox{
	\begin{minipage}{0.5\textwidth}
		More worksheets are available at
		
		\href{https://www.okuyakl.de}{www.okuyakl.de}
	\end{minipage}
	\hfill 
	\begin{minipage}{0.4\textwidth} 
		\includegraphics[width=1.5 cm]{../../viecher/zwe03} 
		% \includegraphics[width=3 cm]{qr221} 
		\includegraphics[width=2 cm]{../../viecher/afanticon1}
		
		\end{minipage}}
		
\end{document}